\chapter{JUSTIFICATIVA}
Na justificativa é necessário fazer uma breve apresentação da temática principal e colocar evidente o porquê da escolha do tema. Pode-se colocar impressões pessoais (sem usar a primeira pessoa do singular) que favorecerem desde a escolha do objeto de pesquisa até questões como se trabalhou, ou trabalha, no local a ser investigado.

Apresentar o problema a ser investigado (lembrando que tem que ser no formato de uma pergunta).

A pesquisa será norteada pelas hipóteses abaixo, sendo que no percurso de investigação elas serão refutadas ou confirmadas:
\begin{uemaenum}
    \item Descrição da Hipótese 1;
    \item Descrição da Hipótese 2;
    \item Descrição da Hipótese 3.
\end{uemaenum}
